\subsection{Addition, Subtraction, and Scalar Multiplication}

\begin{definitionbox}[Entrywise operations]
Let $A=(a_{ij})$ and $B=(b_{ij})$ be matrices of the same size $m\times n$.
Their sum and difference are defined by
\[
A+B=(a_{ij}+b_{ij}),
\qquad
A-B=(a_{ij}-b_{ij}).
\]
For a scalar $\lambda$,
\[
\lambda A=(\lambda a_{ij}).
\]
\end{definitionbox}

\begin{interpretationbox}[Why the sizes must agree]
Each entry of $A$ must have a partner in the same position in $B$. If the shapes
differ, there is no complete position-by-position correspondence, so addition and
subtraction are not defined.
\end{interpretationbox}

\begin{examplebox}[Three entrywise calculations]
Let
\[
A=\begin{pmatrix}2&-1\\0&3\end{pmatrix},
\qquad
B=\begin{pmatrix}5&4\\-2&1\end{pmatrix}.
\]
Then
\[
A+B=
\begin{pmatrix}
2+5&-1+4\\
0+(-2)&3+1
\end{pmatrix}
=
\begin{pmatrix}7&3\\-2&4\end{pmatrix},
\]
\[
A-B=
\begin{pmatrix}
2-5&-1-4\\
0-(-2)&3-1
\end{pmatrix}
=
\begin{pmatrix}-3&-5\\2&2\end{pmatrix},
\]
and
\[
-3A=
\begin{pmatrix}-6&3\\0&-9\end{pmatrix}.
\]
\end{examplebox}

\begin{theorembox}[Rules inherited from ordinary arithmetic]
For matrices of the same size,
\[
A+B=B+A,
\qquad
(A+B)+C=A+(B+C).
\]
More generally, the usual linear rules hold because every statement reduces to
the corresponding statement for each individual entry.
\end{theorembox}

\begin{remarkbox}[Do not transfer this picture to multiplication]
These operations are defined entry by entry. Matrix multiplication will combine
rows with columns instead, so its definition and properties must be developed
separately.
\end{remarkbox}

\begin{summarybox}[Position-by-position arithmetic]
Entrywise operations preserve matrix size and respect fixed positions. Always
check compatibility of shapes before calculating.
\end{summarybox}
